\subsection{Respuesta térmica del marcador y coste de preparación}
\label{sec:termica-respuesta-preparacion-recuperada}

El marcador distingue el estado de referencia dentro del portador completo.
Su probabilidad, su respuesta a una perturbación y el coste de preparar ese
sector son tres lecturas calculables de la misma distribución, con la memoria
conservada. Se recupera así el significado térmico del normalizador, además
de su evaluación numérica.

Consideremos el portador con grados y multiplicidades
\[
 (N,L;d)=(0,0;1),(1,2;380),(2,0;1),(2,1;24),(2,2;624).
\]
Para una separación energética $a>0$, sean $H=aL$, $Q=N-L$ y
$\mu=-a$. El generador de Gibbs es entonces
\[
 K=H-\mu Q=aN,\qquad q=e^{-\beta a},\qquad
 Z_N=\operatorname{Tr}(q^N)=1+380q+649q^2,
 \qquad \rho_N=\frac{q^N}{Z_N}.
\]
Aquí $\beta=(k_B T)^{-1}$, y $P$ es el proyector sobre la referencia marcada
de grado $N=0$. Esa referencia no se confunde con el estado de $L=0$ situado
en el sector $N=2$. Se tiene $p=\operatorname{Tr}(\rho_NP)=Z_N^{-1}$.

\begin{proposition}[Susceptibilidad del sector marcado]
Al perturbar el generador mediante $K_u=K-uP$, manteniendo fijo el portador,
la respuesta $\chi_P=\left.\partial_u\operatorname{Tr}(\rho_uP)\right|_{u=0}$
satisface
\[
 \chi_P=\beta p(1-p),\qquad
 Z_N=1+\frac{\chi_P}{\beta p^2}.
\]
\end{proposition}
\begin{proof}
Como $P$ conmuta con $K$, la perturbación multiplica por $e^{\beta u}$
el peso del sector marcado y deja los demás pesos inalterados. Por tanto,
\[
 p(u)=\frac{p e^{\beta u}}{1-p+p e^{\beta u}}.
\]
La derivación en $u=0$ da $\chi_P=\beta p(1-p)$. Sustituyendo
$p=Z_N^{-1}$ se obtiene la segunda identidad. Equivalentemente,
$P^2=P$ implica $\chi_P=\beta\operatorname{Var}_{\rho_N}(P)$.
\end{proof}

\begin{proposition}[Entropía relativa de la preparación]
Sea $\rho_P=P\rho_NP/p$. Con
$s(\sigma)=-\operatorname{Tr}(\sigma\ln\sigma)$ y
$\mathcal G(\sigma)=\operatorname{Tr}(K\sigma)-\beta^{-1}s(\sigma)$,
\[
 D(\rho_P\Vert\rho_N)=\ln Z_N,\qquad
 \mathcal G(\rho_P)-\mathcal G(\rho_N)=\beta^{-1}\ln Z_N.
\]
\end{proposition}
\begin{proof}
En el soporte de $P$, la relación $\rho_P=\rho_N/p$ da
$\ln\rho_P-\ln\rho_N=-\ln p$, y su traza contra $\rho_P$ es
$-\ln p=\ln Z_N$. Por otra parte,
$\ln\rho_N=-\beta K-\ln Z_N$ implica, para cualquier estado $\sigma$,
\[
 D(\sigma\Vert\rho_N)
 =\beta\bigl[\mathcal G(\sigma)-\mathcal G(\rho_N)\bigr].
\]
La elección $\sigma=\rho_P$ prueba la igualdad energética.
\end{proof}

La diferencia de potencial expresa el trabajo reversible de preparación
cuando se dispone del control y de las reservas térmica y de carga
correspondientes. La identidad describe esa preparación: no atribuye un
consumo fijo a cualquier observación. Si el registro incluye una memoria
independiente $\eta$, la sustitución
$(\rho_N,P)\mapsto(\rho_N\otimes\eta,P\otimes I)$ conserva $p$, la
susceptibilidad y la entropía relativa; el estado $\eta$ permanece intacto.

\subsection{Entropía y capacidad calorífica del portador graduado}
\label{sec:termica-ejemplo-calorifico-recuperado}

En la realización espectral $H_{\rm rel}=3\epsilon N$, la condición
$3\beta\epsilon=\ln1000$ fija $q=10^{-3}$. El normalizador y los dos
momentos del grado son
\[
 Z_N=1{,}380649,\qquad
 \langle N\rangle=\frac{381298}{1380649},\qquad
 \operatorname{Var}(N)=\frac{382842620000}{1906191661201}.
\]
En efecto, las masas de probabilidad de los grados $0,1,2$ son
$1/Z_N$, $380q/Z_N$ y $649q^2/Z_N$. De ellas se obtienen
$\langle N\rangle=(380q+1298q^2)/Z_N$ y
$\langle N^2\rangle=(380q+2596q^2)/Z_N$; su diferencia cuadrática
produce la varianza indicada.

Al conservar el Hamiltoniano y el portador al variar la temperatura,
la entropía $S=k_Bs(\rho_N)$ y la capacidad calorífica $C_V$ resultan
\begin{align*}
 \frac{S}{k_B}
 &=\ln Z_N+\ln1000\,\langle N\rangle
   =2{,}230289295827403\ldots,\\
 \frac{C_V}{k_B}
 &=(\ln1000)^2\operatorname{Var}(N)
   =9{,}58357621855691\ldots.
\end{align*}
La primera fórmula se sigue de $\ln\rho_N=N\ln q-\ln Z_N$.
La segunda utiliza $C_V=k_B\beta^2\operatorname{Var}(H_{\rm rel})$.
Son observables de esta realización térmica. El desplazamiento del origen
espectral registrado por $23$ modifica la energía de preparación y la
intensidad respecto de la referencia; al mantenerse fijo, no altera las
probabilidades condicionadas ni la varianza del grado. La distribución
normalizada y el registro del origen conservan así funciones distintas
dentro del mismo estado.
